% ex: ts=2 sw=2 sts=2 et filetype=tex
% SPDX-License-Identifier: CC-BY-SA-4.0

\documentclass[12pt,addpoints]{exam}

\usepackage[utf8]{inputenc}
\usepackage[T1]{fontenc}
%\usepackage[spanish]{babel}
\usepackage[letterpaper]{geometry}
%\usepackage{graphicx}
%\usepackage{multicol}
%\setlength{\columnsep}{2cm}
%\usepackage{amssymb}

\pagestyle{headandfoot}
\headrule
\header{Pensamiento matemático II}{Examen Parcial III}{CBTIS 246}
\footer{Calificó:~\makebox[4in]{\hrulefill}}{}{Página \thepage\ de \numpages}

\pointpoints{punto}{puntos}
\renewcommand{\solutiontitle}{\noindent\textbf{Solución: }\par\noindent}

%\printanswers

\begin{document}
%\begin{center}
%\fbox{\fbox{\parbox{5.5in}{\centering
%Lee con atención cada pregunta y responde en
%el espacio ubicado en la parte izquierda.
%}}}
%\end{center}
%
%\vspace{5mm}

Nombre:\enspace\hrulefill

\vspace{5mm}

Grupo:\enspace\hrulefill
\enspace{}Grado:\enspace\hrulefill
\enspace{}Fecha:\enspace\hrulefill

\begin{questions}

\begin{EnvFullwidth}
  \sffamily\textbf{Completa la siguiente tabla:}
\end{EnvFullwidth}

% ex: ts=2 sw=2 sts=2 et filetype=tex
% SPDX-License-Identifier: CC-BY-SA-4.0

\question 

\begin{tabular}{|l|c|c|c|}
  \hline
  Por  & $4x^2$  & $-2y^3$  & $5a$ \\ \hline
   & & & \\
  $3x$ & \fillin[$12x^3$] & \fillin[$-6xy^3$] & \fillin[$15xa$] \\ \hline
   & & & \\
  $-y$ & \fillin[$-4x^2y$] & \fillin[$2y^4$] & \fillin[$-5ya$] \\ \hline
   & & & \\
  $2a^4$ & \fillin[$8a^4x^2$] & \fillin[$-4a^4y^3$]  & \fillin[$10a^5$] \\ \hline
\end{tabular}



\begin{EnvFullwidth}
  \sffamily\textbf{Resuelve las siguientes suma de polinomios:}
\end{EnvFullwidth}

% ex: ts=2 sw=2 sts=2 et filetype=tex
% SPDX-License-Identifier: CC-BY-SA-4.0

\question $(3x^2 + 2x + 5) + (2x^2 + 5x + 10)$

  \begin{solutionorlines}[2cm]
    $5x^2 + 7x + 15$
  \end{solutionorlines}

% ex: ts=2 sw=2 sts=2 et filetype=tex
% SPDX-License-Identifier: CC-BY-SA-4.0

\question $(10x^2 + 9x - 4) + (3x^2 - 2x - 7)$

  \begin{solutionorlines}[2cm]
    $13x^2 + 7x - 11$
  \end{solutionorlines}


\begin{EnvFullwidth}
  \sffamily\textbf{Resuelve las siguientes resta de polinomios:}
\end{EnvFullwidth}

% ex: ts=2 sw=2 sts=2 et filetype=tex
% SPDX-License-Identifier: CC-BY-SA-4.0

\question $(11x^3 + 2x^2) - (-3x^3 - 4x^2)$

  \begin{solutionorlines}[2cm]
    $14x^3 + 6x^2$
  \end{solutionorlines}

% ex: ts=2 sw=2 sts=2 et filetype=tex
% SPDX-License-Identifier: CC-BY-SA-4.0

\question $(5a^2 + 2a + 7) - (-4a^2 + 2a - 10)$

  \begin{solutionorlines}[2cm]
    $9a^2 + 17$
  \end{solutionorlines}

% ex: ts=2 sw=2 sts=2 et filetype=tex
% SPDX-License-Identifier: CC-BY-SA-4.0

\question $(4x - 3y) - (10x - 10y)$

  \begin{solutionorlines}[2cm]
    $-6x + 7y$
  \end{solutionorlines}


\begin{EnvFullwidth}
  \sffamily\textbf{Resuelve las siguientes multiplicaciones de polinomios:}
\end{EnvFullwidth}

% ex: ts=2 sw=2 sts=2 et filetype=tex
% SPDX-License-Identifier: CC-BY-SA-4.0

\question $(5x + 3y)(2x - 4y)$

  \begin{solutionorlines}[2cm]
    $10x^2 - 14xy - 12y^2$
  \end{solutionorlines}

% ex: ts=2 sw=2 sts=2 et filetype=tex
% SPDX-License-Identifier: CC-BY-SA-4.0

\question $(-3a^2b + 5a)(-5a + 2a^2b - 10)$

  \begin{solutionorlines}[2cm]
    $25a^3b - 6a^4b^2 + 30a^2b - 25a^2 - 50a$
  \end{solutionorlines}

% ex: ts=2 sw=2 sts=2 et filetype=tex
% SPDX-License-Identifier: CC-BY-SA-4.0

\question $(2x + 4)(2x + 4)$

  \begin{solutionorlines}[2cm]
    $4x^2 + 16x + 16$
  \end{solutionorlines}


\end{questions}

\end{document}
